\cotefile{157-5}
\scanpage{618}
\mspage{6}
alors a) on a $\Phi_1 \mathrel{\longleftarrow\!\mid\!\longrightarrow} \Psi$ et b) $\Phi_1^{\sim}=\Phi^{\sim}=\Phi_1^{\natural}$.

Ce corollaire résulte du théorème en vertu de
\uline{$\overline{\mathrm{XV}}$} prop. 4.2, dont je rappelle ici l'essentiel

\uline{Proposition} 1.3. Soit $\Phi \mathrel{\longleftarrow\!\mid\!\longrightarrow} \Psi$. Alors
$\Phi^{\sim}=\Phi^{\natural}=\Psi^{*}$, et dualement on a $\overline{\Psi}=\Psi^{\natural}=\Phi_{\unsure{*|\star}}$.

Notons aussi, \unsure{dans|en} cet ordre d'idées

\uline{\deleted{Proposition}\inserted{Corollaire} 1.4} Soit $\Phi\longleftrightarrow\Psi$, et soit $W$
une partie \deleted{\unsure{quelconque}}\inserted{quelconque} de $\mathrm{Fl}\,M$. \deleted{\unsure{et W, pour}}

Pour qu'on ait $\Phi\subset W^{\natural}$, il suffit que
toute flèche de $M$ se factorise en $f=pi$,
avec $i\in W$, $p\in\Psi$ \deleted{\unsure{si}}. (Et cette condition
est \inserted{évidemment} aussi nécessaire si $\Phi \mathrel{\longleftarrow\!\mid\!\longrightarrow} \Psi$, et $W$
stable par facteurs directs.)

On dit qu'une partie $\Phi_0$ de $\Phi$ est \deleted{une}
$Q$-\uline{génératrice}\inserted{\unsure{à g.}} de $\Phi$ si elle satisfait les conditions

\deleted{\unsure{a)} $\unsure{\Phi_0}$ \unsure{est} \illegible{1}}

\[
\deleted{\unsure{b)}}\ \Phi\subset\Phi_0^{\sim},\tag*{(1.9)}
\]

et on dit que $\Phi$ est une partie \uline{spéciale} \deleted{\illegible{1}}
si $\Phi$ admet une \inserted{\uuline{petite}} partie $Q$-génératrice à gauche.
Si $\Phi_0\subset\Phi$ est petit, et si $M$ satisfait les
conditions du th. 1.1, alors $\Phi_0$, est une
partie $Q$-génératrice (à gauche) sss son
$Q$-saturé total (à g.) $\Phi_0'''$ contient $\Phi$
(condition en apparence plus forte que $\Phi\subset\Phi_0^{\sim}$,
puisque en général on a $\Phi_0'''\subset\Phi_0^{\sim}$). Cela
signifie donc \unsure{aussi|que}, via 1.2, que \deleted{\unsure{toute}}
% archive sheet number: 617
\scanpage{619}
\mspage{7}
élément de $\Phi$ est facteur direct d'un composé
transfini de flèches déduites de $\Phi_0\cup(\text{Identités de }M)$
par cochgt. de base.

Je rappelle encore les sorites $\pm$ tautologiques
des foncteurs adjoints

\uline{Proposition} 1.5. Soit

\diagram{\begin{tikzcd}[ampersand replacement=\&,column sep=large]
\text{(1.5.1)} \& M' \arrow[r,shift right,"u^{\unsure{\alpha|*}}"'] \& M \arrow[l,shift right,"u_{!}"']
\end{tikzcd}}

un couple de foncteurs adjoints, \unsure{Soient|soient} \deleted{$\Phi\unsure{\in|\subset}\mathrm{Fl}\,M$},
\deleted{partie $\mathrm{Fl}\,M$, $\Psi'$ une \unsure{part}}

\[
\Phi\subset\mathrm{Fl}(M),\qquad \Psi'\subset\mathrm{Fl}\,M'), 
\]

alors on a

\[
\big(u_!(\Phi)\longleftrightarrow\Psi'\big)\Longleftrightarrow
\big(\Phi\longleftrightarrow u^*(\Psi')\big)\tag*{(1.5.2)}
\]

\uline{Corollaire 1.6}. Pour \unsure{tte|toute} partie $\Phi$ de $\mathrm{Fl}(M)$,
\deleted{\unsure{posant} $\Psi=\Phi_*$, \unsure{on}}

\[
\big(u_!(\Phi)\big)_*=(u^*)^{-1}(\Phi_*)\tag*{(1.6.1)}
\]

En effet, appliquant 1.5. à $\Phi$ et à $\{f'\}$,
où $f'\in\mathrm{Fl}\,M$, on trouve

\[
f'\in\big(u_!(\Phi)\unsure{\big\rangle|\big)}_*\Longleftrightarrow u^*(f)\in\Phi_*,
\]

ce qui équivaut à (1.6.1). Dualement

\uline{Corollaire 1.7} : Pour tte partie $\Psi'$ de $M'$, on
a

\[
\big(u^*(\Psi')\big)^*=(u_!)^{-1}(\Psi'{}^*)\tag*{(1.7.1)}
\]

% slanted
\margin{\unsure{donner ici} 1,9. \unsure{pag|page} 9}

En effet, si $f\in\mathrm{Fl}\,M$, on a par (1.5.2)

\[
f\in\unsure{(|\{}u^*(\Psi')\unsure{\}|)}^{*}\Longleftrightarrow u_!(f)\in\Psi'{}^*,
\]
% archive sheet number: 618
\scanpage{620}
\mspage{8}
Un couple $Q$-orthogonal \inserted{$\Phi\longleftrightarrow\Psi$} est dit \deleted{\unsure{$Q$|s}} \uline{spécial}
s'il satisfait les conditions suivantes

\[
\left\{\begin{aligned}
\text{a)}\quad &M\ \text{est \unsure{stable|stble} par \inserted{petites} }\unsure{\underline{\lim}|\varinjlim},\ \text{les objets de }M\text{ sont}\\
&\text{accessibles}\quad [\text{On dit alors que }M\text{ est ``\unsure{spécial|spéciale}''}]\\
\text{b)}\quad &\deleted{\Phi\ \text{\unsure{est}}}\ \text{Le couple est clos : }\Phi\longleftrightarrow\Psi.\\
\text{c)}\quad &\Phi\ \text{est \uline{spéciale}.}
\end{aligned}\right.\tag*{(1.10)}
\]

Cela implique donc, par 1.1, que $\Phi \mathrel{\longleftarrow\!\mid\!\longrightarrow} \Psi$ est
un couple de Quillen. Si $M$ satisfait a),
les couples $Q$-orthogonaux spéciaux dans $M$,
ou encore les couples de Quillen spéciaux,
correspondent donc biunivoquement aux
parties $\Phi$ de $\mathrm{Fl}\,M$ qui sont \uline{spéciales} et
closes à gauche (ou ce qui revient au $=$, $Q$-saturées
à gauche), en \deleted{\unsure{fai}} associant à une telle
$\Phi$ le couple $\Phi \mathrel{\longleftarrow\!\mid\!\longrightarrow} \Phi_{\unsure{*|\alpha}}$.

Ceci posé :

\uline{Corollaire} 1.8 : Supposons que $M,M'$ satisfassent
\margin{\tikz[x=1em,y=1em]{\draw (0,0) .. controls (-.08,-1.3) and (.12,-2.6) .. (.22,-3.6); \draw (.065,.06) .. controls (.035,-1.4) and (.25,-2.5) .. (.12,-3.45);}}
la condition a) de (1.10), et soit $(\Phi,\Psi)$
un couple de Quillen spécial dans $M$. Alors
$\big(\Phi'=u_!(\Phi)^{\sim},\ \Psi'=(u^*)^{-1}(\deleted{\unsure{M}}\Psi)\big)$ est un
couple de Quillen spécial dans $M'$.

En effet, par le cor. 1.6 on a

\[
\Phi'_*=\big(u_!(\Phi)\big)_{\unsure{*|\alpha}}=(u^{*})^{-1}(\Psi)=\Psi',
\]

donc $\Phi'\longleftrightarrow\Psi'$. Il reste à prouver que
$\Phi'$ est \unsure{spécial|spéciale}. Soit $\Phi_0$ une partie \unsure{petite} de
$\Phi$ telle que $\Phi\subset\Phi_0^{\sim}$, posons $\Phi_0'=u_!(\Phi_0)\inserted{\subset}\Phi'$
% archive sheet number: 619
\scanpage{621}
\mspage{9}
il suffit de voir qu'on a

\[
u_!(\Phi)\subset u_!(\Phi_0)^{\sim}\tag*{$(*)$}
\]

d'où résultera qu'on a $\Phi'\overset{\mathrm{def}}{=}\big(u_!(\Phi)\big)^{\sim}\subset u_!(\Phi_0)^{\sim}$.
Mais on a

\[
u_!(\Phi_0)^{\sim}=\big(u_!(\Phi_0)_*\big)^*\quad \deleted{\text{et }u_!(\Phi_0}
\]

et

\[
u_!(\Phi_0)_*=(u^{\unsure{*|\alpha}})^{-1}(\underbrace{\Phi_{0*}}_{\Psi})=u^{*\, -1}(\Psi)\overset{\mathrm{def}}{=}\Psi',
\]

donc

\[
u_!(\Phi_0)^{\sim}=\Psi'{}^*=\Phi'=u_!(\Phi)^{\sim}
\]

d'où \unsure{(*)}. Cet argument ne fait pas appel
à l'hypothèse ``spéciale'' sur $M$, et peut
s'expliciter en le

% slanted
\margin{devrait passer avant 1.8.}
\uline{Corollaire} 1.9. Sous les conditions générales de
\deleted{\unsure{(}}1.5\deleted{\unsure{)}} ($u_!:M\deleted{\unsure{^{\natural}}}\to M'$ étant un foncteur admettant
un adjoint à droite)\deleted{\unsure{$u^*$}}, soient $\Phi_0,\Phi$ deux
parties de $\mathrm{Fl}\,M$ telles que $\Phi\subset\Phi_0^{\sim}$ (ou ce
qui revient au $=$, $\Phi^{\sim}\subset\Phi_0^{\sim}$). Alors

\[
u_!(\Phi)\subset u_!(\Phi_0)^{\sim}\qquad\text{i.e.}\quad u_!(\Phi)^{\sim}\subset u_!(\Phi_0)^{\sim}.
\]

En effet, cette relation équivaut aussi à
$u_!(\Phi_0)_*\subset u_!(\Phi)_*$, c'est-à-dire (par 1.6)
$(u^*)^{-1}(\Phi_{0*})\subset(u^*)^{-1}(\Phi_*)$, et est \unsure{impliqué|impliquée}
par $\Phi_{0*}\subset\Phi_*$. Cette dernière \unsure{d'autre part}
équivaut à l'hypothèse $\Phi^{\sim}\subset\Phi_0^{\sim}$, \unsure{qed}.

\unsure{!}
% archive sheet number: 620
\scanpage{622}
\mspage{10}
\deleted{\unsure{Soit}} Si $\Psi$ est une partie de $\mathrm{Fl}\,M$, on en déduit
pour \unsure{tte} $I$ dans $\mathrm{Cat}$ une partie

\[
\left\{\begin{aligned}
&\Psi(I\unsure{\rangle|)}\subset\mathrm{Fl}\,M(I),\qquad\text{avec }M(I)=\underline{\mathrm{Hom}}(I^{\mathrm o},M),\\
&f\in\Psi(I)\ \Longleftrightarrow\ f(i)\in\Psi\quad \forall i\in\mathrm{Ob}\,I.
\end{aligned}\right.\tag*{(1.11)}
\]

\deleted{\unsure{Soit}} Si on introduit la catégorie discrète $I_0$
ss-jacente à $\mathrm{Ob}\,I$, et le foncteur canonique

\[
\unsure{v|u}_I=v:I_0\longrightarrow I,\tag*{(1.12)}
\]

on trouve un couple de foncteurs adjoints

\diagram{\begin{tikzcd}[ampersand replacement=\&,row sep=0pt,column sep=large]
\text{(1.1\deleted{2}3)} \& M(I)\arrow[r,shift right=1.2,"v^*"'] \& M(I_0)\simeq M^{I_0}\arrow[l,shift right=1.2,"v_!"',"{\unsure{\cdot|>}}"{description,pos=.08}] ,\\
\& {[}\phantom{M(I)} \& \phantom{M(I_0)}{]}\arrow[l,dotted,"v_*"]
\end{tikzcd}}

et on a

\[
\Psi(I)=(v^*)^{-1}\big(\Psi(I_0)\big)\unsure{;|.}
\]

\deleted{donc par 1.6}
\[
\deleted{\Psi(I)=v_!(\Phi(I_0))}
\]

Supposons maintenant qu'on ait

\[
\left\{\begin{aligned}
&\Psi=\Phi_*\\
&\text{d'où}\qquad\Psi(I_0)=\Phi(I_0)_*
\end{aligned}\right.\tag*{(1.14)}
\]

(vérification immédiate), on conclut alors par 1.6

\[
\Psi(I)=\big(v_!(\Phi(I_0)\big)_*=\big(\Phi\langle I\rangle\big)_*\tag*{(1.15)}
\]

avec

\[
\Phi\langle I\rangle\overset{\mathrm{def}}{=}v_!\big(\Phi(I_0)\big)^{\sim}=\Psi(I)^*.\tag*{(1.16)}
\]

Si on a un foncteur

\[
u:I\longrightarrow I',\tag*{(1.17)}
\]

d'où un triple de foncteurs adjoints,
% archive sheet number: 621
\scanpage{623}
\mspage{11}
\diagram{\begin{tikzcd}[ampersand replacement=\&,column sep=large]
\text{(1.18)} \& \left\{\vphantom{\dfrac{M}{M}}\right. M(I) \arrow[r,shift left=3,"u_!"] \arrow[r,shift right=3,"u_*","{\deleted{\unsure{u_*}}}"'] \& M(I')\arrow[l,"u^*"' ] ,
\end{tikzcd}}

on aura de façon $\pm$ tautologique

\[
\left\{\begin{aligned}
&u_!\big(\Phi\langle I\rangle\big)\subset\Phi\langle I'\rangle\\
&u^*\big(\Psi(I')\big)\subset\Psi(I).
\end{aligned}\right.\tag*{(1.19)}
\]

Le corollaire 1.8. donne : \deleted{\unsure{donc}}

\uline{Proposition} 1.10 : Supposons que \deleted{$\Phi,\Psi$} $\Phi \mathrel{\longleftarrow\!\mid\!\longrightarrow} \Psi$
soit un couple de Quillen \uline{spécial} dans $M$,
\unsure{i.e.} de la forme $(\Phi,\Phi_*)$, avec $\Phi$ partie
$Q$-close à gauche spéciale dans $M$ ($M$
satisfaisant la condition (1.10. a)). Alors
pour \unsure{tt} $I$ dans $\mathrm{Cat}$, $(\Phi\langle I\rangle,\Psi(I))$
est un couple spécial dans $M(I)$, en
particulier il satisfait au th. de factorisation.

Considérons la construction duale,
qui nous fait in\deleted{\unsure{tro}}troduire, à partir d'un
couple $\Phi\longleftrightarrow\Psi$ dans $M$, \unsure{des|les}
couples

\[
\tikz[baseline=(rel.base)]{
  \node[inner xsep=.45em,inner ysep=.6ex] (rel)
    {$\Phi(I\unsure{\rangle|)}\longleftrightarrow\Psi[I]$};
  \draw ([xshift=-1.3em,yshift=-.3ex]rel.south west)
    -- ([xshift=.25em]rel.south west) -- ([xshift=-.23em]rel.north west)
    -- ([yshift=.45ex]rel.north east) -- (rel.south east)
    -- ([xshift=.5em]rel.south east);
},\tag*{(1.20)}
\]

où $\Phi(I)$ est défini par la formule générale
(1.11) (où $\Psi$ est remplacé par $\Phi$), et où
$\Psi[I]$ est défini soit comme $\Phi(I)_*$,
soit comme

\[
\Psi[I]=\unsure{v|u}_*\big(\Psi(I_0)\big)^{-}=\Phi(I)_*.\tag*{(1.2\unsure{1|4})}
\]
% archive sheet number: 622
\scanpage{624}
\mspage{\uline{12}}
On a alors, pour $u:I\longrightarrow I'$

\[
\left\{\begin{aligned}
&u^*(\Phi(I'))\subset\deleted{\unsure{\Phi}}\Phi(I)\\
&u_*(\Psi[I])\subset\Psi[I']
\end{aligned}\right.\tag*{(1.22)}
\]

Il se pose encore la question si le couple
$(\Phi(I),\Psi[I])$ satisfait au \deleted{le} th. de factorisation,
sous des hypothèses convenables sur $(\Phi,\Psi)$.
Je peux seulement hasarder la

\uline{Question-conjecture} \uline{1.11}. Supposons que
\deleted{$(\Phi,\unsure{\Psi}$} $(\Phi,\Psi)$ soit un couple de Quillen
\uline{spécial}, en est-il de $=$ de $(\Phi(I),\Psi[I])$
pour \unsure{tt} $I$ dans $(\mathrm{Cat})$ ? \deleted{\unsure{Je co}} En d'autres
termes, si \deleted{$\Phi$ \unsure{est une}} $M$ est spécial
(condition (1.10 a)), et si $\Phi$ est une
partie $Q$-close à g. \uline{spéciale} de $M$ i.e.
admettant une \deleted{petite} partie $Q$-génératrice,
est-il vrai que pour \unsure{tte} $I$ dans $\mathrm{Cat}$,
$\Phi(I)$ est également spéciale ?
\inserted{Je conjecture que oui, en faisant au \unsure{besoin} des hypothèses}
% slanted
\margin{\unsure{suppl.} sur $M$ (existence d'une filtration \unsure{``cardinale''} ?)}
P. ex. si $\Phi_0$
est $Q$-génératrice dans $\Phi$, et \deleted{petite}, alors
$\Phi_0(I)$ est-elle $Q$-génératrice \unsure{dans} $\Phi(I)$ ?

\uline{Prop. 1.12} Soit $\Phi\longleftrightarrow\Psi$ dans $M$, soit $u:I\longrightarrow I'$
une flèche de $\mathrm{Cat}$, d'où

\diagram{\begin{tikzcd}[ampersand replacement=\&,column sep=large]
\text{(1.12.1)} \& M(I)\arrow[r,shift left=2.5,"u_!"]\arrow[r,shift right=2.5,"u_*"'] \& M(I')\arrow[l,"u^*"] .
\end{tikzcd}}

% slanted
\margin{cf \uline{$\overline{\mathrm{XVII}}$} 6.4 (p. \unsure{117|217|17}), mais où les notations pour $\Phi\langle I\rangle,\Psi(I)$ sont inversées.}
On a alors (1.19) \unsure{sc.|sic.|i.e.}

\[
\left\{\begin{aligned}
&u_!\Phi(\langle I\rangle)\subset\Phi\langle I'\rangle\\
&u^*(\Psi(\unsure{I|I'}))\subset\Psi(I'),
\end{aligned}\right.\tag*{(1.13.1)}
\]

et de plus l'équivalence

\[
u^*(\Phi\langle I'\rangle)\subset\Phi\langle I\rangle\Longleftrightarrow
u_*(\Psi(I))\subset\Psi(I').\tag*{(1.13.2)}
\]

Ces relations équiv. sont satisfaites dans chacun des 3 cas :
$1^\circ)$ $u$ est coétale
$2^\circ)$ $I$ est discrète
$3^\circ)$ $\forall i'$ dans $I'$, \deleted{toute} \inserted{\unsure{chaq.}} \unsure{comp. conn.} de $I_{/i'}$ admet un \unsure{objet} final.
% slanted
\margin{\uline{NB} On a $1^\circ\Rightarrow3^\circ$, $2^\circ\Rightarrow3^\circ$}
% archive sheet number: 623
\scanpage{625}
\mspage{13}
2. \uline{Rappels sur les triples de Quillen} \inserted{clos}. \uline{\deleted{\unsure{sur les}} Triples de
Quillen spéciaux}.

Soit $(M,W)$ une \uline{catégorie de \unsure{modèle|modèles}}, i.e
$M$ muni d'un \uline{localiseur} $W\subset\mathrm{Fl}(M)$, i.e une
partie satisfaisant

\[
\begin{aligned}
&W\text{ contient les \unsure{iso|isos}, et si }f,g\text{ sont des}\\
&\text{flèches composables, si deux parmi }f,g,gf\\
&\text{sont dans }W,\text{ la troisième aussi}\\
&\text{(\uline{saturation faible}).}
\end{aligned}\tag*{(2.1)}
\]

Je rappelle que la donnée d'un \uline{triple de
Quillen clos}\inserted{$(W,C,F)$}, ou d'une structure de Quillen
close, sur $M$, ayant $W$ comme ensemble des
quasi-isomorphismes, revient à la donnée de
deux couples de Quillen

\[
C \mathrel{\longleftarrow\!\mid\!\longrightarrow} TF,\qquad TC \mathrel{\longleftarrow\!\mid\!\longrightarrow} F,\tag*{(2.2)}
\]

satisfaisant

\[
TC\subset C\cap W,\qquad TF\subset F\cap W,\tag*{(2.3)}
\]

plus une des trois conditions équivalentes

\[
\left\{\begin{aligned}
\text{a)}\quad &TC=C\cap W &(\text{i.e. } C\cap W\subset TC)\\
\text{b)}\quad &TF=\deleted{\unsure{C}}F\cap\deleted{\unsure{F}}W &(\text{i.e. }\deleted{\unsure{C}}F\cap W\subset T\deleted{\unsure{C}}F)\\
\text{c)}\quad &W=TF\circ TC
\end{aligned}\right.\tag*{(2.4)}
\]

Rappelons aussi que ceci implique que $W$ \uline{est fortement saturée},
et a fortiori que $W=W^{\natural}$ i.e
$W$ stable par facteurs directs.

\unsure{Rappellons|Rappelons} aussi qu'en vertu de (1.4), \deleted{les} \inserted{on a}
\deleted{conditions 2.3 \unsure{a) b)} équivalent respectivement à}
\inserted{ceci :}
% archive sheet number: 624
\scanpage{626}
\mspage{\deleted{\unsure{13}} \unsure{(14)|14}}
\deleted{(2.5) a') Toute flèche $f$ de $M$ (ou seulement, toute
flèche de $C\cap W$) \unsure{se} factorise en $pi$, avec
$i$}

(2.5) \unsure{Ètant|Étant} donnés les deux couples (2.2), avec $TC\subset C$,
\deleted{$TF\subset F$} (équivaut à $TF\subset F$), pour qu'on ait
la première relation (2.3) i.e $TC\subset W$, il \deleted{suffit, si}
\unsure{faut}, et il suffit si $W=W^{\natural}$, que \unsure{tte|toute} flèche $f$
de $M$ (ou seulement \unsure{tte} flèche de $TC$) se
factorise en $f=pi$, avec $i\in W$, $p\in F$. Dualement,
pour qu'on ait $TF\subset W$, il faut, et il
suffit si $W=W^{\natural}$, que \unsure{tte} flèche $f$ de $M$
(ou seulement \unsure{tte|toute} flèche de $T\unsure{F|C}$) se factorise
en $f=pi$, avec $i=W$, $f\in C$.

\deleted{\uline{Proposition 2.1} Soit}
\diagram{\begin{tikzcd}[ampersand replacement=\&,column sep=large,execute at end picture={\draw (current bounding box.south west) -- (current bounding box.north east);}]
\deleted{M}\arrow[r,shift left,"{\deleted{\unsure{u|v}_!}}"] \& \deleted{M'}\arrow[l,shift left,"{\deleted{u^*}}"]
\end{tikzcd}}
\deleted{un couple de foncteurs adjoints, et \unsure{soit}}

On dit que la structure de Quillen $(W,C,F)$
est \uline{spéciale} si les deux couples de Quillen
(2.3) sont spéciaux, i.e. si les conditions
suivantes sont satisfaites

\[
\begin{aligned}
\text{a)}\quad &M\ \text{est stable par \inserted{petites} }\unsure{\underline{\lim}|\varinjlim},\ \text{et ses éléments}\\
&\text{sont accessibles}\qquad(M\text{ ``spécial''}).\\
\text{b)}\quad &\text{Il existe des \uline{petites} parties }C_0\subset C,\\
&TC_0\subset TC\ \text{qui soient } Q\text{-génératrices à g.,}\\
&\text{i.e. telles que l'on ait}\\[-1pt]
&\left\{\begin{aligned}
&C\subset C_0^{\sim} &&\text{i.e. } \unsure{C|\subset}=C_0^{\sim} &&\text{i.e. } TF=C_{0*}\\
&TC\subset(TC_0)^{\sim} &&\text{i.e. } TC=(TC_0)^{\sim} &&\text{i.e } F=(TC_0)_*
\end{aligned}\right.
\end{aligned}\tag*{(2.6)}
\]

Ainsi, la donnée d'un triple de Quillen
\uline{spécial} dans $M$, \deleted{\unsure{satis}} quand $M$ satisfait (2.6 a),
revient à la donnée des deux parties \deleted{\unsure{closes}}
% archive sheet number: 625
\scanpage{627}
\mspage{15}
\[
TC\subset C\tag*{(2.7)}
\]
de $\mathrm{Fl}\,M$, satisfaisant les conditions

% slanted
\margin{$C,TC$ ``spéciales''}
\[
\left|\begin{aligned}
&\left\{\begin{aligned}
\text{a)}\quad &\deleted{\unsure{\Phi}}C,TC\text{ sont }Q\text{-closes à gauche}\\
\text{b)}\quad &\deleted{\unsure{\Phi}}C,TC\text{ sont }Q\text{-engendrées à g. par}\\
&\text{des \uline{petites} parties } C_0,TC_0\text{ de }C,TC\\
&\text{respectivement}
\end{aligned}\right.\\
&\text{c)}\quad TC\subset W,\quad TF\overset{\mathrm{def}}{=}C_*\subset W\\
&\text{d)}\quad\text{On a l'une des \inserted{trois} conditions équivalentes}\\
&\qquad\text{supplémentaires de (2.4), i.e}\\
&\qquad\alpha)\ TC=C\cap W\quad\text{i.e}\quad C\cap W\subset TC\\
&\qquad\beta)\ TF=F\cap W\quad\text{i.e.}\quad F\cap W\subset TF\\
&\qquad\gamma)\ W=TF\circ TC.
\end{aligned}\right.\tag*{(2.8)}
\]

\uline{Proposition} 2.1. Soit

\diagram{\begin{tikzcd}[ampersand replacement=\&,column sep=large]
\text{(2.1.1)} \& M'\arrow[r,shift right,"u^*"'] \& M\arrow[l,shift right,"u_!"']
\end{tikzcd}}

un couple de foncteurs adjoints, $M$ étant
muni d'une structure de Quillen close
$(W,C,F)$. Posons

\[
\left\{\begin{aligned}
&\deleted{C'=u_!(C)^{\sim}}\\
&\deleted{TC'=u_!(TC)^{\sim}}\\
&\deleted{F'\unsure{.}}
\end{aligned}\right.\quad
\left\{\begin{aligned}
&W'=(u^*)^{-1}(W)\\
&F'=(u^*)^{-1}(W)\\
&TF'=F'\cap W'=(u^*)^{-1}(TF)\\
&\deleted{\unsure{T}}C'=u_!(C)^{\sim}\\
&TC'=u_!(TC)^{\sim}.
\end{aligned}\right.\tag*{(2.1.2)}
\]

$1^\circ)$ On a alors

\[
\left\{\begin{aligned}
&C'\longleftrightarrow TF',\qquad TC'\longleftrightarrow F'\\
&T\unsure{\cdot} C'\subset C',\qquad TF'\subset F',
\end{aligned}\right.\tag*{(2.1.3)}
\]

Pour que les données (2.1.\unsure{1|2}) dans $M'$ \unsure{//|))} correspondent à un triple de Quillen $(W',C',F')$, il
faut et il suffit que les conditions \unsure{suivantes|équiv. suivantes}
soient satisfaites :
% archive sheet number: 626
\scanpage{628}
\mspage{16}
\[
\left\{\begin{aligned}
\text{a)}\quad &\text{Factorisation pour } C'\longleftrightarrow TF'\text{ et }TC'\longleftrightarrow F',\\
&\text{i.e. les relations } C' \mathrel{\longleftarrow\!\mid\!\longrightarrow} TF'\text{ et }TC' \mathrel{\longleftarrow\!\mid\!\longrightarrow} F'.\\
\text{b)}\quad &TC'\subset W'.
\end{aligned}\right.\tag*{(2.1.4)}
\]

ou encore, au lieu de b), la condition équivalente
(en présence de a)) :

\[
\left\{\begin{aligned}
\text{b')}\quad &\text{Toute flèche }f'\text{ dans }M'\text{ (ou seulement}\\
&f'\in TC')\text{ se factorise en }f'=p'i',\text{ avec}\\
&i'\in W',\ p'\in F'.
\end{aligned}\right.\tag*{(2.1.5)}
\]

$2^\circ)$ Supposons \inserted{(2.1.4.} a)) satisfaite \inserted{(factorisation)}. Alors, pour que
$(W',C',F')$ soit un triple de Quillen clos,
il \uline{suffit} que l'on ait

\[
\text{b'')}\quad\left|\underline{u^*u_!(TC)\subset TC}\right|,\tag*{(2.1.6)}
\]

et que de plus on soit dans l'un des
trois cas suivants

A) $u^*$ admet un adjoint à droite $u_*$

B) $u^*$ commute aux \deleted{\unsure{pet}}\unsure{petites} $\varinjlim$, $(W,C,F)$
est spécial, et $M'$ est spécial i.e. satisfait à la
condition (2.6 a).

C) $u^*$ est pleinement fidèle, et $u^*u_!(F)\subset F$.

Notons tout de suite le

\uline{Corollaire} 2.2. Supposons que $(W,C,F)$ soit un
\deleted{\unsure{triple}} de Quillen spécial, que $M'$ soit spécial,
\deleted{et qu'on ait} que $u^*$ commute aux petites $\varinjlim$,
enfin que l'on ait (2.1.6) i.e. $u^*u_!(TC)\subset TC$.
Alors $(W',C',F')$ est un couple de Quillen
spécial.
% archive sheet number: 627
\scanpage{629}
\mspage{17}
En effet, en vertu de 2.1. $2^\circ$, cas B, \inserted{pour \deleted{\unsure{voir}} que $(W',C',F')$ soit un triple de Quillen clos,} il reste
à vérifier \deleted{que} (2.1.4\deleted{\unsure{b)}}\unsure{a)}). Mais cela résulte
de 1.8. appliqué aux deux couples de Quillen
spéciaux $(C,TF),(TC,F)$ dans $M$ : les couples
$(C',TC')$ et $(TC',F')$ sont des couples de Quillen
spéciaux, ce qui prouve en même temps
que le triple $(W',C',F')$ est spécial, \unsure{qed}.

Démontrons 2.1. \deleted{\unsure{Les}} les \inserted{premières deux} relations (2.1.3) sont
contenues dans 1.6, les relations $TC'\subset C'$, $TF'\subset F'$
sont tautologiques. Le critère (2.1.4) pour
que $(W',C',F')$ soit un triple de Quillen clos,
est contenu dans le rappel de la page 13, compte
tenu que l'on sait déjà que $TF'=F'\cap W'$
(relation (2.4 b)). Que la condition (1.1.4 b),
en présence de (2.1.4. a), soit équivalente
à la condition (2.1.5 b'), est contenu dans
la première assertion (2.5), \unsure{en notant|où notant} que
l'on a $W'=W'{}^{\natural}$ (ce qui résulte de la définition
de $W'$ et de $W^{\natural}=W^{\natural}$).

Il reste à prouver \inserted{$2^\circ)$, i.e} que la condition b) est
\deleted{\unsure{la forme équi}} satisfaite \deleted{\unsure{dès}} moyennant \inserted{a) et} b'') (2.1.6),
dans chacun des trois cas A), B), C).

\uline{Cas A} Par définition de $W'=(u^*)^{-1}(W)$, pour
prouver que \deleted{\unsure{u}} $TC'\subset W'$, il suffit de prouver

\[
\deleted{\unsure{(2.1.7)}}\qquad u^*(TC')\subset W'
\]

et a fortiori, il suffit de prouver

\[
u^*\deleted{\unsure{u}}(TC')\subset\deleted{\unsure{W}}TC\quad\text{i.e.}\quad u^*(TC')\longleftrightarrow F\tag*{(2.1.7)}
\]

Or par 1.5. appliqué au couple de foncteurs
% archive sheet number: 628
\scanpage{630}
\mspage{18}
adjoints $(u^*,u_*)$, cette relation équivaut à

\[
TC'\unsure{\longleftrightarrow|\Longleftrightarrow}u_*(F).
\]

Comme $TC'\overset{\mathrm{def}}{=}u_!(TC)^{\sim}$, ceci équivaut à

\[
u_!(TC)\longleftrightarrow u_*(\unsure{F}),
\]

ou encore (en réappliquant 1.5 à $u_!(TC),F$
au lieu de $(TC',F)$), à

\[
u^*u_!(TC)\longleftrightarrow F\qquad\text{i.e.}\quad
u^*u_!(TC)\subset\underset{\substack{\parallel\\ F^*}}{TC},
\]

or c'est là l'hypothèse \deleted{\unsure{a)}} b'').

\uline{Cas B} Il suffit encore de prouver
2.1.7. Or, du fait que $u_!(TC)$ \deleted{\unsure{admet}}\inserted{\unsure{admet}} \unsure{est}
\unsure{un} \tikz[baseline=(under.base)]{\node[inner sep=0pt] (under) {\deleted{\unsure{petit ensemble} $Q$-générateur}}; \draw[dashed] ([xshift=.6em,yshift=-.2ex]under.south west) -- ([xshift=-.3em,yshift=-.2ex]under.south east);} à gauche,
$u_!(TC)^{\sim}$ est la clôture de $u_!(TC)$ pour les
opérations de cochgt. de base, de composition
transfinie, de facteur direct. Comme
$u^*$ commute à \unsure{ttes|toutes} ces opérations (comme
il commute aux \inserted{petites} $\unsure{\underline{\lim}|\varinjlim}$), et \deleted{\unsure{étant donn}} comme
$TC$ est dans $\mathrm{Fl}\,M$ est stable par ces opérations,
on voit donc que pour montrer qu'\unsure{on|en}
$u^{\unsure{*|'}}$ applique $u_!(TC)^{\sim}$ dans $TC$, il suffit de
voir qu'il applique $u_!(TC)$ dans $TC$, ce qui
est encore l'hypothèse b'').

\uline{Cas C} On va vérifier b) sous la forme
équivalente b'). Soit donc $f'\in\mathrm{Fl}\,M'$,
\deleted{\unsure{factorisons} \unsure{u} \unsure{soit}} $f':X'\unsure{\text{---}|\longrightarrow}Y^{\unsure{F|'|*}}$,
% archive sheet number: 629
\scanpage{631}
\mspage{19}
\deleted{factorisons} posons

\[
(\qquad f=u_{\deleted{\unsure{!|*}}}^{*}(f'),
\]

et factorisons $f$ en $pi$, avec $i\in TC$, $p\in F$,
d'où

\[
u_!u^{\unsure{*|u}}(f')=u_{\deleted{\unsure{*}}!}(f)=u_!(p)\,u_!(i).
\]

Mais $u^*$ étant \unsure{pl.} fidèle $u_!u^*\xrightarrow{\sim}\mathrm{id}$, donc
on trouve une factorisation de \deleted{$\unsure{u_!(f)}$}
$f'$ en $u_!(p)\,u_!(i)$. \deleted{Comme $u_!(F)\subset F'$} Il suffit
donc de voir que

\[
u_!(p)\in F',\qquad u_!(i)\in W',
\]

ce qui équivaut à

\[
u^*u_!(p)\in F,\qquad u^*u_!(i)\in W
\]

et résulte des hypothèses \deleted{\unsure{b}} que $u^*u_!(F)\subset F$
et $u^*u_!(TC)\subset TC\subset W$, qed.

\uline{NB} J'ai noté les cas A) et C) par acquit de
conscience, sans application \unsure{précise|pratique} en vue. Le cas
\inserted{le plus} utile sera \deleted{\unsure{le}}\inserted{\unsure{justiciable du}} cas B), \deleted{qui}\inserted{et} mérite d'être
\unsure{réconforté|récapitulé} en un

% slanted
\margin{cf Scholie 2.10, qui devrait précéder le th. 2.3.}
\uline{Théorème} 2.3. Soit $(W,C,F)$ un triple de Quillen
spécial \unsure{dans} $M$. Pour \unsure{tt} $I$ dans $\mathrm{Cat}$, considérons
le triple $(W(I),C\langle I\rangle,F(I))$\inserted{\deleted{\unsure{dans}} $M(I)=\underline{\mathrm{Hom}}(I^{\mathrm o},M)$}, où les notations
sont celles des pages 10, donc $W(I),F(I)$ et
$TF(I)=\deleted{\unsure{T}}F(I)\cap W(I)$ se définissant fibre par fibre,

\[
C\langle I\rangle=(F(I))_{\deleted{\unsure{*}}}^{*}=\big(\unsure{v|u}_{I!}(C(I_0))\big)^{\sim}.
\]
% slanted
\margin{où $v_I:I_0\to I$ ($I_0$ cat. discrète \inserted{ayant \unsure{$=$|mêmes} objets} \deleted{\unsure{de cat.} \illegible{1}} \unsure{que} $I$) \deleted{\unsure{source de $I$}} et}
\deleted{Alors}

a) Pour \unsure{tte} $I$ dans $\mathrm{Cat}$, $(W(I),C\langle I\rangle,F(I))$ est
% archive sheet number: 630
\scanpage{632}
\mspage{20}
un triple de Quillen spécial dans $M(I)$.

$2^\circ)$ Si $u\unsure{:|\in} I\unsure{\text{---}|\longrightarrow} I'$ est une flèche dans $\mathrm{Cat}$, d'où

\diagram{\begin{tikzcd}[ampersand replacement=\&,column sep=large]
\text{\uline{(2.3.1)}} \& M(I)\arrow[r,shift left,"u_!"] \& M(I')\arrow[l,shift left,"u^*"]
\end{tikzcd}}

\[
\left\{\begin{aligned}
u^*(W(I'))&\subset W(I)\\
u^*(F(I'))&\subset F(I)
  &\smash{\tikz[baseline=(note.base),x=1em,y=1em]{
    \node[anchor=base west,inner sep=0pt] (note) at (.7,0) {tautologiques};
    \draw (0,1.82) .. controls (-.06,1.30) and (-.17,.49) .. (.13,.27)
      .. controls (-.03,.06) and (.25,-.50) .. (.18,-.72)
      .. controls (.17,-.91) and (.06,-.99) .. (0,-1.04);
  }}\\
u^*(TF(\unsure{I'|J'}))&\subset TF(I)\\
u_!(\unsure{C|TC}\langle I\rangle)&\subset C\langle I'\rangle\\
u_!(TC\langle I\rangle)&\subset TC\langle I'\rangle,
\end{aligned}\right.\tag*{(2.3.2)}
\]

enfin, $u_!$ induit un foncteur

\[
M(I)_c\longrightarrow M(I')_c\tag*{(2.3.3)}
\]

\margin{\begin{tabular}{|@{\,}l}sous-catégories\\ des objets\\ cofibrants\end{tabular}}
\uline{compatible aux localiseurs} $W(I)_c=W(I)\mid M(I)_c$
et $W(I')_c=W(I')\mid M(I')_c$, donc induit
un foncteur

\[
u_!^W:\mathrm{Ho}_W(I)\longrightarrow\mathrm{Ho}_W(I'),\tag*{(2.3.\unsure{4|5})}
\]

Ce foncteur est un adjoint à gauche du
foncteur $u_W^{\unsure{*}}$,

\[
u_{\unsure{W|*W}}^{\deleted{\unsure{W|*}}}:\mathrm{Ho}_W(I')\longrightarrow\mathrm{Ho}_W(I)\tag*{(2.3.5)}
\]

induit par $u^*$ \deleted{\unsure{de}} (2.3.1).

\uline{Démonstration de} 2.3. Prouvons \unsure{$1^\circ$)}. On considère
comme dans \unsure{1.10.|n. 10.|p. 10.} le foncteur

\[
\unsure{v|u}=v_I:I_0\unsure{\text{---}|\longrightarrow} I
\]

et le couple de foncteurs adjoints

\diagram{\begin{tikzcd}[ampersand replacement=\&,column sep=large]
\text{(2.3.6)} \& M(I)\arrow[r,shift right,"v^*"'] \& M(I_0)\arrow[l,shift right,"v_!"'] ,
\end{tikzcd}}

auquel on peut appliquer 2.2, \deleted{\unsure{cas B}} Il reste
% archive sheet number: 631
